\begin{examplebox}
\textbf{\textcolor{CExample}{例 1（直接求积：识别信号即用）}}
\par\textbf{题目：}椭圆 $\dfrac{x^2}{4}+\dfrac{y^2}{3}=1$，左右顶点
$A_1(-2,0),A_2(2,0)$。已知 $P(1,\frac32)$ 在椭圆上，求 $k_{PA_1}k_{PA_2}$。
\par\textbf{解：}核验 $\frac14+\frac{(3/2)^2}{3}=\frac14+\frac34=1$。
$a^2=4,b^2=3$，故
\[
 \boxed{k_{PA_1}k_{PA_2}=-\frac{b^2}{a^2}=-\frac34}.
\]
\textbf{避坑：}点属于椭圆是使用第三定义的前提；这里恰好满足。

\medskip
\textbf{\textcolor{CExample}{例 2（逆向求椭圆、离心率）}}
\par\textbf{题目：}椭圆 $C:\dfrac{x^2}{a^2}+\dfrac{y^2}{b^2}=1\ (a>b>0)$，
$P(2,1)\in C$，且 $k_{PA_1}k_{PA_2}=-\frac12$，求 $C$ 的方程及离心率。
\par\textbf{解：}由第三定义，$b^2/a^2=1/2$，即 $b^2=a^2/2$。
代入点 $P$，
\[
 \frac4{a^2}+\frac1{b^2}=1
 \ \Longrightarrow\ \frac4{a^2}+\frac2{a^2}=1
 \ \Longrightarrow\ a^2=6,\ b^2=3.
\]
故
\[
 \boxed{C:\frac{x^2}{6}+\frac{y^2}{3}=1},\qquad
 \boxed{e=\sqrt{1-\frac{b^2}{a^2}}=\frac{\sqrt2}{2}}.
\]
\textbf{更快：}斜率积 $K=e^2-1$，于是直接有 $e=\sqrt{1-\frac12}=\frac{\sqrt2}{2}$。

\medskip
\textbf{\textcolor{CExample}{例 3（高价值综合：直径圆过两个定点）}}
\par\textbf{题目：}椭圆 $\dfrac{x^2}{4}+\dfrac{y^2}{3}=1$ 的左右顶点
$A_1(-2,0),A_2(2,0)$，$P\ne A_1,A_2$ 在椭圆上。
$PA_1$ 与 $x=4$ 交于 $M$，$PA_2$ 与 $x=4$ 交于 $N$。
证明以 $MN$ 为直径的圆恒过两个定点，并求其坐标。
\begin{center}
\begin{tikzpicture}[x=.85cm,y=.85cm,font=\small]
 \path[use as bounding box] (-3,-4.4) rectangle (8,4.2);
 \node at (2.5,3.85) {两条顶点连线截出 $MN$，直径圆固定穿过 $S,T$};
 \draw[->,gray] (-2.7,0)--(7.7,0) node[right] {$x$};
 \draw[blue!60!black,thick] (0,0) ellipse (2 and {sqrt(3)});
 \draw[orange!80!black,thick] (4,0) circle (3);
 \draw[gray,dashed] (4,-3.4)--(4,3.4);
 \node[right,gray] at (4,-.8) {$x=4$};
 \draw[teal!70!black,thick] (-2,0)--(4,3);
 \draw[magenta!70!black,thick] (1,1.5)--(4,-3);
 \draw[orange!80!black,very thick] (4,3)--(4,-3);
 \fill (-2,0) circle (1.5pt) node[below left] {$A_1$};
 \fill (2,0) circle (1.5pt) node[above right] {$A_2$};
 \fill (1,1.5) circle (1.5pt);
 \node[above] at (1,2) {$P$};
 \fill (4,3) circle (1.5pt) node[above right] {$M$};
 \fill (4,-3) circle (1.5pt) node[below right] {$N$};
 \fill (1,0) circle (1.5pt) node[below left] {$S(1,0)$};
 \fill (7,0) circle (1.5pt) node[below right] {$T(7,0)$};
 \node at (2.5,-4.1) {图取 $P=(1,\frac32)$；一般情形仍有 $y_My_N=-9$};
\end{tikzpicture}
\end{center}

\par\textbf{解：}记 $k_1=k_{PA_1},k_2=k_{PA_2}$，
由两条直线的点斜式
\[
 y=k_1(x+2),\quad y=k_2(x-2)
 \quad\Longrightarrow\quad M=(4,6k_1),\ N=(4,2k_2).
\]
由第三定义 $k_1k_2=-\frac34$，故
\[
 y_My_N=(6k_1)(2k_2)=12\Bigl(-\frac34\Bigr)=-9.
\]
以 $MN$ 为直径的圆满足
\[
 (x-4)^2+(y-y_M)(y-y_N)=0.
\]
令 $y=0$，上式变为 $(x-4)^2+y_My_N=0$，
即 $(x-4)^2-9=0$，从而 $x=1$ 或 $x=7$。
\[
 \boxed{\text{两个定点为 }(1,0),\ (7,0).}
\]
\textbf{突破口：}直径圆上固定点是否存在，不必先求圆心和半径；优先考虑
$y=0$，可避开易变的 $y_M+y_N$，只留下已知定值 $y_My_N$。

\medskip
\textbf{\textcolor{CExample}{例 4（长轴竖直：现场推导不混淆）}}
\par\textbf{题目：}$\dfrac{x^2}{4}+\dfrac{y^2}{9}=1$，上下顶点
$B_1(0,-3),B_2(0,3)$。取短轴端点 $P(2,0)$，求 $k_{PB_1}k_{PB_2}$。
\par\textbf{解：}因长轴竖直，$a=3,b=2$，
\[
 \boxed{k_{PB_1}k_{PB_2}=-\frac{a^2}{b^2}=-\frac94}.
\]
直接检查：$k_{PB_1}=\frac{0+3}{2}=\frac32$，
$k_{PB_2}=\frac{0-3}{2}=-\frac32$，其积确为 $-\frac94$。

\medskip
\textbf{\textcolor{CExample}{例 5（逆向轨迹与判断长短轴）}}
\par\textbf{题目：}固定 $A_1(-3,0),A_2(3,0)$，动点 $P$ 满足
$k_{PA_1}k_{PA_2}=-\frac49$。求轨迹方程及离心率。
\par\textbf{解：}设 $P(x,y)$，则
\[
 \frac{y^2}{x^2-9}=-\frac49
 \ \Longleftrightarrow\ \frac{x^2}{9}+\frac{y^2}{4}=1.
\]
因 $4<9$，$A_1,A_2$ 确为长轴端点，$a=3,b=2$，故
\[
 \boxed{e=\frac{\sqrt{a^2-b^2}}a=\frac{\sqrt5}{3}}.
\]
\textbf{严格轨迹：}应写作椭圆 $\frac{x^2}{9}+\frac{y^2}{4}=1$
\textbf{去掉} $(-3,0),(3,0)$，因为斜率在这两点不成立。

\medskip
\textbf{\textcolor{CExample}{例 6（从圆压到椭圆，再从椭圆拉回圆：逐点核算）}}
\par\textbf{题目：}设 $a=2,\ b=\frac65$（即 $b=1.2$），圆
$C_0:u^2+v^2=4$ 上有点 $Q=(\frac85,\frac65)$。
令 $A_1(-2,0),A_2(2,0)$。按保持横坐标的纵向缩放得到点 $P$。
求 $P$ 的坐标，验证它在椭圆上，并分别计算压缩前后的斜率乘积。
\par\textbf{第 1 步（核查圆点）：}
$\left(\frac85\right)^2+\left(\frac65\right)^2=\frac{64+36}{25}=4$，所以 $Q\in C_0$。
\par\textbf{第 2 步（坐标映射）：}压缩系数是 $\frac ba=\frac35$；横坐标不改，纵坐标乘 $\frac35$，得到
\[
 \boxed{P=\left(\frac85,\frac35\cdot\frac65\right)
 =\left(\frac85,\frac{18}{25}\right)}.
\]
\par\textbf{第 3 步（验证椭圆）：}
\[
 \frac{(8/5)^2}{4}+\frac{(18/25)^2}{36/25}
 =\frac{16}{25}+\frac9{25}=1.
\]
即 $P$ 在 $\frac{x^2}{4}+\frac{y^2}{36/25}=1$ 上。
\par\textbf{第 4 步（斜率究竟变了多少）：}
\[
 \begin{aligned}
  k_{QA_1}&=\frac{6/5}{8/5+2}=\frac13,&
  k_{QA_2}&=\frac{6/5}{8/5-2}=-3,\\
  k_{PA_1}&=\frac{18/25}{8/5+2}=\frac15=\frac35k_{QA_1},&
  k_{PA_2}&=\frac{18/25}{8/5-2}=-\frac95=\frac35k_{QA_2}.
 \end{aligned}
\]
因此圆内斜率积是 $\frac13\cdot(-3)=-1$，
椭圆内斜率积是 $\frac15\cdot(-\frac95)=\boxed{-\frac9{25}}$
$=-\left(\frac35\right)^2$。
\par\textbf{第 5 步（反向还原）：}把 $P$ 的纵坐标乘 $\frac ab=\frac53$，
$\frac53\cdot\frac{18}{25}=\frac65$，恰好恢复原点 $Q=(\frac85,\frac65)$。
\textbf{迁移要点：}端点与横向位移不变；纵向增量变为原来的 $3/5$，
两条斜率也各变为原来的 $3/5$。\textbf{不是}每条连线的长度都乘 $3/5$。

\medskip
\textbf{\textcolor{CExample}{高价值变式（先思考再看答案）}}
\par\textbf{题目：}长轴水平的椭圆上，$k_{PA_1}k_{PA_2}=-\frac{7}{16}$，
不求椭圆方程，直接求 $e$。若斜率积为 $-\frac{16}{7}$，是否还能维持“长轴水平”的前提？
\par\textbf{答案：}$e^2=1-\frac{7}{16}=\frac9{16}$，故
$\boxed{e=\frac34}$。后一种\textbf{不可能}：长轴水平时 $0<b^2/a^2<1$，
所以 $-1<k_{PA_1}k_{PA_2}<0$，不可能小于 $-1$。
\end{examplebox}
